% =====================
% Caso lineal
% =====================
\caption{Campo incidente lineal en base cartesiana. Se muestran $|E_x|^2$, $|E_y|^2$ e intensidad total.}
\label{fig:lineal-incidente-componentes}

\caption{Campo en el plano focal para incidencia lineal. La componente $|E_y'|^2$ corresponde a la polarización cruzada.}
\label{fig:lineal-plano-focal-componentes}

\caption{Mapa de polarización del campo incidente lineal. El fondo muestra $|E|^2$ y las flechas la orientación local de la polarización.}
\label{fig:lineal-incidente-mapa-polarizacion}

\caption{Mapa de polarización en el plano focal para incidencia lineal. El fondo muestra $|E|^2$ y las flechas la orientación local, evidenciando la rotación inducida por la componente cruzada.}
\label{fig:lineal-plano-focal-mapa-polarizacion}

% =====================
% Caso circular
% =====================
\caption{Campo incidente circular en base $(L,R)$. Se muestran $|E_L|^2$, $|E_R|^2$ e intensidad total.}
\label{fig:circular-incidente-componentes}

\caption{Campo en el plano focal para incidencia circular. La componente opuesta representa la conversión parcial de helicidad.}
\label{fig:circular-plano-focal-componentes}

\caption{Mapa de polarización del campo incidente circular. El fondo muestra $|E|^2$ y las elipses el estado local de polarización.}
\label{fig:circular-incidente-mapa-polarizacion}

\caption{Mapa de polarización en el plano focal para incidencia circular. El fondo muestra $|E|^2$ y las elipses la polarización local, incluyendo la modificación espacial de la helicidad.}
\label{fig:circular-plano-focal-mapa-polarizacion}

% =====================
% Mapas escalares de polarización
% =====================
\caption{Mapa escalar del ángulo de elipticidad $\chi$ en el plano focal para incidencia lineal. Las regiones en blanco corresponden a zonas de intensidad demasiado baja, donde el parámetro deja de estar bien definido.}
\label{fig:lineal-plano-focal-elipticidad}

\caption{Mapa escalar del ángulo del eje mayor $\psi$ en el plano focal para incidencia lineal. Las regiones en blanco corresponden a zonas de intensidad demasiado baja, donde el parámetro deja de estar bien definido.}
\label{fig:lineal-plano-focal-psi}

\caption{Mapa escalar del ángulo de elipticidad $\chi$ en el plano focal para incidencia circular. Las regiones en blanco corresponden a zonas de intensidad demasiado baja, donde el parámetro deja de estar bien definido.}
\label{fig:circular-plano-focal-elipticidad}

\caption{Mapa escalar del ángulo del eje mayor $\psi$ en el plano focal para incidencia circular. Las regiones en blanco corresponden a zonas de intensidad demasiado baja, donde el parámetro deja de estar bien definido.}
\label{fig:circular-plano-focal-psi}

% =====================
% Caso cilíndrico
% =====================
\caption{Campo incidente radial en base cilíndrica. Se muestran $|E_\rho|^2$, $|E_\varphi|^2$ e intensidad total.}
\label{fig:radial-incidente-componentes}

\caption{Campo en el plano focal para incidencia radial, representado en base cilíndrica.}
\label{fig:radial-plano-focal-componentes}

\caption{Mapa de polarización del campo incidente radial. El centro se excluye por la singularidad de la base cilíndrica.}
\label{fig:radial-incidente-mapa-polarizacion}

\caption{Mapa de polarización en el plano focal para incidencia radial, con muestreo polar y centro excluido.}
\label{fig:radial-plano-focal-mapa-polarizacion}

\caption{Campo incidente azimutal en base cilíndrica. Se muestran $|E_\rho|^2$, $|E_\varphi|^2$ e intensidad total.}
\label{fig:azimutal-incidente-componentes}

\caption{Campo en el plano focal para incidencia azimutal, representado en base cilíndrica.}
\label{fig:azimutal-plano-focal-componentes}

\caption{Mapa de polarización del campo incidente azimutal. El centro se excluye por la singularidad de la base cilíndrica.}
\label{fig:azimutal-incidente-mapa-polarizacion}

\caption{Mapa de polarización en el plano focal para incidencia azimutal, con muestreo polar y centro excluido.}
\label{fig:azimutal-plano-focal-mapa-polarizacion}
